% Generated draft: requires mathematical and source review.
\documentclass[11pt]{article}
\usepackage[margin=25mm]{geometry}
\usepackage{fontspec}
\setmainfont{Latin Modern Roman}
\usepackage{amsmath,amssymb,mathtools,unicode-math}
\setmathfont{Latin Modern Math}
\usepackage{enumitem,needspace,xurl,hyperref}
\setlength{\parindent}{0pt}
\setlength{\parskip}{5pt}
\setlength{\emergencystretch}{3em}
\newcommand{\N}{\mathbb{N}}
\newcommand{\Z}{\mathbb{Z}}
\newcommand{\Q}{\mathbb{Q}}
\newcommand{\R}{\mathbb{R}}
\newcommand{\C}{\mathbb{C}}
\newcommand{\F}{\mathbb{F}}
\newcommand{\A}{\mathbb{A}}
\newcommand{\PP}{\mathbb{P}}
\newcommand{\E}{\mathbb{E}}
\newcommand{\eps}{\varepsilon}
\newcommand{\dd}{\,\mathrm{d}}
\begin{document}
\section*{UnsolvedMath ID: 30000382}
\textbf{Source problem code:} catalogue entry
\subsection{Definitions and mathematical statement}
\textbf{Bedrock uniqueness question.} If a universe \(V\) satisfies the Bedrock Axiom, is its bedrock model unique? In the usual ZFC setting, a ground of \(V\) is a transitive inner model \(W\subseteq V\), containing all ordinals and satisfying ZFC, for which there are set forcing \(\mathbb P\in W\) and a \(W\)-generic filter \(G\subseteq\mathbb P\) with \(V=W[G]\). The trivial extension is permitted, so \(V\) may itself be a ground. Class forcing is not substituted for set forcing.\par The Ground Axiom for a model \(W\) says that \(W\) is not a set-forcing extension of any proper inner ground. A bedrock is a ground \(W\) of \(V\) satisfying the Ground Axiom, described in the source as a minimal ground. The Bedrock Axiom asserts that at least one such ground exists. The requested uniqueness statement is\[\forall W_0,W_1\,\bigl[(W_0,W_1\text{ are grounds of }V\text{ satisfying the Ground Axiom})\ \Longrightarrow\ W_0=W_1\bigr].\] The equality is equality of the inner models, not just isomorphism, and the question concerns uniqueness conditional on existence, not whether every universe has a bedrock. A minimal ground should not silently be defined to be a least ground contained in all others, which would build a stronger conclusion into the terminology.\par The supplied literature triage checked 2026-08-21 labels the question open and says surveys continue to list it among basic unresolved problems. It supplies a survey reference but no exact later theorem statements or verification of their scope. This is a dated supplied status assertion; the mathematical question and definitions are retained without independently asserting that bedrock uniqueness remains unresolved today.
\subsection{Short English statement}
Assuming a set-forcing ground with no deeper proper ground exists, can there be two different such grounds? The problem asks for uniqueness of the inner model, conditional on existence. Minimal is not silently replaced by least, and the supplied open-status assessment is unverified dated curation.
\subsection{Sources}
\begin{itemize}\item Local UnsolvedMath registry, permanent ID 30000382, code catalogue entry; \emph{Uniqueness of Bedrock Models}, Oberwolfach Reports Open Problems.\item \emph{Set Theory (2006)}, Oberwolfach Reports, DOI 10.4171/owr/2005/55, pp. 3121--3174. \url{https://doi.org/10.4171/owr/2005/55}. Proposed year and year tag: 2005; no proposer named.\item Supplied literature label \emph{original Oberwolfach report}: \url{https://doi.org/10.4171/owr/2005/55}.\item Supplied literature label \emph{set theoretic geology survey}: \url{https://doi.org/10.1016/j.apal.2014.11.004}.\item Literature triage checked 2026-08-21, registry status \texttt{open}; record updated 2026-09-16. The \texttt{corrected\_verified} field records notation restoration only. Verify the dated unresolved-status claim independently.\end{itemize}
\end{document}
